\documentclass[12pt,a4paper]{article}

% 中文支持 - 使用Windows标准字体（SimSun/SimHei）
\usepackage{xeCJK}
\setCJKmainfont{SimSun}      % 宋体（正文）
\setCJKsansfont{SimHei}      % 黑体（标题）
\setCJKmonofont{Courier New}  % 等宽字体（代码）

% 定义中文等宽字体族（用于代码块）
\setCJKfamilyfont{zhfs}{Courier New}  % 中文等宽字体

% 定义黑体字体族
\setCJKfamilyfont{zhhei}{SimHei}  % 黑体字体族定义

% 页面设置（A4，2.5cm边距）
\usepackage[top=2.5cm, bottom=2.5cm, left=2cm, right=2cm]{geometry}

% 标题格式
\usepackage{titlesec}
% 一级标题：黑体18pt，段前12pt，段后6pt
\titleformat{\section}
  {\fontsize{18pt}{18pt}\bfseries\CJKfamily{zhhei}}
  {}
  {0em}
  {}
\titlespacing{\section}{0pt}{12pt}{6pt}

% 二级标题：黑体16pt，段前6pt，段后3pt
\titleformat{\subsection}
  {\fontsize{16pt}{16pt}\bfseries\CJKfamily{zhhei}}
  {}
  {0em}
  {}
\titlespacing{\subsection}{0pt}{6pt}{3pt}

% 三级标题：黑体14pt，段前3pt，段后3pt
\titleformat{\subsubsection}
  {\fontsize{14pt}{14pt}\bfseries\CJKfamily{zhhei}}
  {}
  {0em}
  {}
\titlespacing{\subsubsection}{0pt}{3pt}{3pt}

% 四级标题：黑体12pt，段前3pt，段后0pt
\titleformat{\paragraph}
  {\fontsize{12pt}{12pt}\bfseries\CJKfamily{zhhei}}
  {}
  {0em}
  {}
\titlespacing{\paragraph}{0pt}{3pt}{0pt}

% 段落格式
\usepackage{setspace}
\onehalfspacing  % 1.5倍行距
\setlength{\parindent}{2em}  % 首行缩进2字符
\setlength{\parskip}{6pt}    % 段后间距6pt

% 颜色支持（必须在listings之前加载）
\usepackage{xcolor}    % 颜色支持

% 代码块格式
\usepackage{listings}  % 代码块支持

% 代码块格式设置
\lstset{
    basicstyle=\ttfamily\CJKfamily{zhfs}\fontsize{10.5pt}{12.6pt}\selectfont,  % 代码字体：Courier New + 中文，五号（10.5pt）
    breaklines=true,              % 自动换行
    breakatwhitespace=true,      % 在空格处换行
    frame=single,                % 边框
    framesep=3pt,                % 边框间距
    framerule=0.5pt,             % 边框宽度
    rulecolor=\color{black},    % 边框颜色
    backgroundcolor=\color{gray!5},  % 背景色（浅灰色）
    showstringspaces=false,      % 不显示字符串中的空格
    showspaces=false,            % 不显示空格
    showtabs=false,              % 不显示制表符
    tabsize=4,                   % 制表符大小
    captionpos=b,                % 标题位置（底部）
    numbers=left,                % 行号位置
    numberstyle=\tiny\color{gray},  % 行号样式
    stepnumber=1,                % 行号步长
    numbersep=5pt,               % 行号与代码间距
    extendedchars=true,          % 支持扩展字符（包括中文）
    inputencoding=utf8           % UTF-8编码支持
}

% 代码块标题格式
\renewcommand{\lstlistingname}{代码}


% 表格格式
\usepackage{tabularx}
\usepackage{longtable}  % 长表格支持
\usepackage{booktabs}
\usepackage{array}
\usepackage{colortbl}  % 表格背景色支持（必须在xcolor之后加载）

% 表格格式设置
\definecolor{tableheaderbg}{RGB}{240,240,240}  % 表头背景色RGB(240,240,240)
\renewcommand{\arraystretch}{1.2}              % 表格行高
\setlength{\arrayrulewidth}{0.5pt}            % 表格内边框宽度

% 表格字体设置（表格文字：宋体，五号10.5pt）
\newcolumntype{C}[1]{>{\centering\arraybackslash\fontsize{10.5pt}{12.6pt}\selectfont}p{#1}}
\newcolumntype{L}[1]{>{\raggedright\arraybackslash\fontsize{10.5pt}{12.6pt}\selectfont}p{#1}}
\newcolumntype{R}[1]{>{\raggedleft\arraybackslash\fontsize{10.5pt}{12.6pt}\selectfont}p{#1}}

% 表格默认宽度设置（占页面宽度的90%-100%）
\setlength{\tabcolsep}{6pt}  % 表格列间距

% 页眉页脚
\usepackage{fancyhdr}
\setlength{\headheight}{14pt}  % 修复页眉高度警告
\pagestyle{fancy}
\fancyhead[L]{\small\bfseries 文档名称}
\fancyhead[R]{\small\bfseries 版本号：v1.0}
\fancyfoot[L]{\small 编制单位}
\fancyfoot[C]{\small 第\thepage\ 页 共\pageref{LastPage}\ 页}
\fancyfoot[R]{\small \today}
\renewcommand{\headrulewidth}{0.5pt}
\renewcommand{\footrulewidth}{0pt}

% 图表编号格式
\usepackage{caption}  % 图表标题格式
\captionsetup{
    labelsep=space,           % 标签和标题之间用空格分隔
    font=small,               % 字体大小
    format=hang,              % 标题格式
    justification=centering,   % 居中
    singlelinecheck=false     % 允许多行标题
}

% 图表编号格式：图1-1、表2-3
\renewcommand{\figurename}{图}
\renewcommand{\tablename}{表}
\renewcommand{\thefigure}{\arabic{section}-\arabic{figure}}
\renewcommand{\thetable}{\arabic{section}-\arabic{table}}

% 表格标题字体：黑体，五号（10.5pt），加粗
\captionsetup[table]{font={bf,small},labelfont=bf}
\captionsetup[figure]{font={bf,small},labelfont=bf}

% 其他设置
\usepackage{amsmath}     % 数学公式支持（numberwithin命令需要此包）
\usepackage{hyperref}
\usepackage{lastpage}
% 公式编号格式：(1-1)、(2-3)
\numberwithin{equation}{section}  % 公式按章节编号
\renewcommand{\theequation}{\arabic{section}-\arabic{equation}}


\begin{document}
\section{GB/T 8567-2006 文档计划标准模板}
\subsection{[列表] 文档基本信息}
\textbf{文档编号}: [项目编号]-DP-001  
\textbf{文档名称}: 文档计划 (Documentation Plan)  
\textbf{项目名称}: [项目名称]  
\textbf{编制单位}: [编制单位名称]  
\textbf{编制日期}: YYYY-MM-DD  
\textbf{版本号}: v1.0  

\hrule
\subsection{文档变更记录}
\begin{longtable}{|p{2.5cm}|p{3cm}|p{3cm}|p{3cm}|p{2.5cm}|}
\hline
版本号 & 变更日期 & 变更内容 & 变更人 & 审核人 \\
\hline
v1.0 & YYYY-MM-DD & 初始版本 & [姓名] & [姓名] \\
\hline
\end{longtable}

\hrule
<!-- 目录会自动生成，使用Pandoc转换时会自动生成目录 -->
<!-- 如果使用Markdown编辑器，可以保留手动目录作为预览 -->

\hrule
\subsection{引言}
\subsubsection{目的}
本文档旨在：
\begin{itemize}
\item 明确项目文档编制的范围、责任和进度安排
\item 规定文档的格式、风格和质量标准
\item 建立文档管理、评审和版本控制机制
\item 确保项目文档的完整性、一致性和可追溯性
\end{itemize}

\subsubsection{范围}
本文档适用于：
\begin{itemize}
\item [项目名称] 项目的所有文档编制工作
\item 项目文档的全生命周期管理
\item 参与项目文档编制的所有人员
\end{itemize}

本文档不适用于：
\begin{itemize}
\item 外部供应商提供的文档
\item 临时性文档和草稿
\end{itemize}

\subsubsection{术语和定义}
\begin{longtable}{|p{4cm}|p{5cm}|p{5cm}|}
\hline
术语 & 英文 & 定义 \\
\hline
文档计划 & Documentation Plan & 规定项目文档编制范围、责任、进度和标准的计划文档 \\
文档编制 & Documentation Development & 按照标准要求编写、评审和发布文档的过程 \\
文档评审 & Documentation Review & 对文档内容、格式和质量进行检查和评估的过程 \\
文档版本 & Document Version & 文档的版本标识，用于跟踪文档的变更历史 \\
\hline
\end{longtable}

\subsubsection{引用文件}
\begin{itemize}
\item GB/T 8567-2006《计算机软件文档编制规范》
\item GB/T 9385-2008《计算机软件需求规格说明规范》
\item GB/T 9386-2008《计算机软件测试文档编制规范》
\item [其他相关标准或文件]
\end{itemize}

\subsubsection{概述}
本文档是项目文档编制的总体规划，依据GB/T 8567-2006《计算机软件文档编制规范》第5.3节"文档计划"的要求编制。

文档计划包括以下主要内容：
\begin{enumerate}
\item 文档编制范围：明确需要编制的文档类型和数量
\item 文档编制责任：明确各类文档的编制、审核和批准责任
\item 文档编制进度：制定文档编制的时间表和里程碑
\item 文档格式和标准：规定文档的格式、风格和质量要求
\item 文档评审与批准：建立文档评审和批准流程
\item 文档版本控制：建立文档版本管理机制
\item 文档存储与管理：建立文档存储和管理机制
\item 文档质量保证：建立文档质量保证机制
\end{enumerate}

\hrule
\subsection{文档编制范围}
\subsubsection{文档列表}
根据GB/T 8567-2006标准和项目实际情况，确定需要编制的文档如下：

\begin{longtable}{|p{2.29cm}|p{2.29cm}|p{2.29cm}|p{2.29cm}|p{2.29cm}|p{2.29cm}|p{2.29cm}|}
\hline
序号 & 文档名称 & 文档编号 & 文档类型 & 优先级 & 编制方式 & 备注 \\
\hline
1 & 可行性研究报告 & FSR & 项目管理 & 高 & 新编制 & - \\
2 & 项目开发计划 & PDP & 项目管理 & 高 & 完善现有 & - \\
3 & 软件需求说明书 & SRS & 需求分析 & 高 & 完善现有 & - \\
4 & 数据要求说明书 & DRS & 需求分析 & 中 & 新编制 & - \\
5 & 概要设计说明书 & HLD & 设计 & 中 & 完善现有 & - \\
6 & 详细设计说明书 & DLD & 设计 & 中 & 新编制 & - \\
7 & 数据库设计说明书 & DBD & 设计 & - & 不适用 & 本项目不涉及数据库 \\
8 & 模块开发卷宗 & MDD & 开发 & 低 & 新编制 & 可选 \\
9 & 测试计划 & TP & 测试 & 中 & 完善现有 & - \\
10 & 测试分析报告 & TAR & 测试 & 中 & 完善现有 & - \\
11 & 用户手册 & UM & 用户文档 & 低 & 完善现有 & - \\
12 & 操作手册 & OM & 用户文档 & 低 & 完善现有 & - \\
13 & 开发进度月报 & MPR & 项目管理 & 低 & 完善现有 & 可选 \\
14 & 项目开发总结报告 & PSR & 项目管理 & 低 & 完善现有 & - \\
\hline
\end{longtable}

\textbf{说明}：
\begin{itemize}
\item 文档编号：按照GB/T 8567-2006标准的文档编号规则
\item 优先级：高（必须编制）、中（建议编制）、低（可选编制）
\item 编制方式：新编制、完善现有、不适用
\end{itemize}

\subsubsection{文档分类}
按文档类型分类：

\textbf{2.2.1 项目管理文档}
\begin{itemize}
\item 可行性研究报告（FSR）
\item 项目开发计划（PDP）
\item 开发进度月报（MPR）
\item 项目开发总结报告（PSR）
\end{itemize}

\textbf{2.2.2 需求分析文档}
\begin{itemize}
\item 软件需求说明书（SRS）
\item 数据要求说明书（DRS）
\end{itemize}

\textbf{2.2.3 设计文档}
\begin{itemize}
\item 概要设计说明书（HLD）
\item 详细设计说明书（DLD）
\item 数据库设计说明书（DBD）- 不适用
\end{itemize}

\textbf{2.2.4 开发文档}
\begin{itemize}
\item 模块开发卷宗（MDD）- 可选
\end{itemize}

\textbf{2.2.5 测试文档}
\begin{itemize}
\item 测试计划（TP）
\item 测试分析报告（TAR）
\end{itemize}

\textbf{2.2.6 用户文档}
\begin{itemize}
\item 用户手册（UM）
\item 操作手册（OM）
\end{itemize}

\subsubsection{文档适用性说明}
\textbf{2.3.1 必须编制的文档（11个）}
上述文档列表中优先级为"高"和"中"的文档，必须按照计划编制完成。

\textbf{2.3.2 可选编制的文档（2个）}
\begin{itemize}
\item 模块开发卷宗（MDD）：根据项目规模和复杂度决定是否编制
\item 开发进度月报（MPR）：根据项目管理需要决定是否编制
\end{itemize}

\textbf{2.3.3 不适用文档（1个）}
\begin{itemize}
\item 数据库设计说明书（DBD）：本项目不涉及数据库设计，故不适用
\end{itemize}

\hrule
\subsection{文档编制责任}
\subsubsection{组织架构}
文档编制组织架构如下：

\begin{verbatim}
项目文档编制组
├── 文档负责人
│   ├── 负责文档计划制定
│   ├── 负责进度跟踪
│   └── 负责质量监督
├── 技术负责人
│   ├── 负责技术文档审核
│   └── 负责技术内容把关
├── 测试负责人
│   ├── 负责测试文档编制
│   └── 负责测试文档审核
└── 文档编写人员
    ├── 负责具体文档编写
    └── 负责文档维护
\end{verbatim}

\subsubsection{职责分工}
\begin{longtable}{|p{3cm}|p{4cm}|p{4cm}|p{3cm}|}
\hline
角色 & 姓名 & 主要职责 & 负责文档 \\
\hline
文档负责人 & [姓名] & 文档计划制定、进度跟踪、质量监督 & 全部文档 \\
技术负责人 & [姓名] & 技术文档审核、技术内容把关 & FSR, PDP, SRS, DRS, HLD, DLD \\
测试负责人 & [姓名] & 测试文档编制和审核 & TP, TAR \\
文档编写人员A & [姓名] & 项目管理文档编写 & FSR, PDP, MPR, PSR \\
文档编写人员B & [姓名] & 需求和分析文档编写 & SRS, DRS \\
文档编写人员C & [姓名] & 设计文档编写 & HLD, DLD \\
文档编写人员D & [姓名] & 测试文档编写 & TP, TAR \\
文档编写人员E & [姓名] & 用户文档编写 & UM, OM \\
\hline
\end{longtable}

\subsubsection{人员配置}
\textbf{3.3.1 人员数量}
\begin{itemize}
\item 文档负责人：1人
\item 技术负责人：1人
\item 测试负责人：1人
\item 文档编写人员：3-5人
\end{itemize}

\textbf{3.3.2 能力要求}
\begin{itemize}
\item 文档负责人：熟悉GB/T 8567-2006标准，具备文档管理经验
\item 技术负责人：熟悉项目技术架构，具备技术文档审核能力
\item 测试负责人：熟悉测试流程，具备测试文档编写能力
\item 文档编写人员：熟悉文档编写规范，具备良好的文字表达能力
\end{itemize}

\hrule
\subsection{文档编制进度}
\subsubsection{总体进度计划}
\textbf{项目文档编制总工期}：12周（3个月）

\textbf{阶段划分}：
\begin{itemize}
\item 阶段1：启动文档（第1-2周）
\item 阶段2：需求文档（第3-4周）
\item 阶段3：设计文档（第5-7周）
\item 阶段4：测试文档（第8-9周）
\item 阶段5：用户文档（第10周）
\item 阶段6：项目管理文档（第11-12周）
\end{itemize}

\subsubsection{详细进度表}
\begin{longtable}{|p{2.67cm}|p{2.67cm}|p{2.67cm}|p{2.67cm}|p{2.67cm}|p{2.67cm}|}
\hline
文档编号 & 文档名称 & 计划开始 & 计划完成 & 持续时间 & 负责人 \\
\hline
FSR & 可行性研究报告 & 第1周 & 第2周 & 2周 & [姓名] \\
PDP & 项目开发计划 & 第1周 & 第2周 & 2周 & [姓名] \\
SRS & 软件需求说明书 & 第3周 & 第4周 & 2周 & [姓名] \\
DRS & 数据要求说明书 & 第3周 & 第4周 & 2周 & [姓名] \\
HLD & 概要设计说明书 & 第5周 & 第6周 & 2周 & [姓名] \\
DLD & 详细设计说明书 & 第5周 & 第7周 & 3周 & [姓名] \\
TP & 测试计划 & 第8周 & 第9周 & 2周 & [姓名] \\
TAR & 测试分析报告 & 第8周 & 第9周 & 2周 & [姓名] \\
UM & 用户手册 & 第10周 & 第10周 & 1周 & [姓名] \\
OM & 操作手册 & 第10周 & 第10周 & 1周 & [姓名] \\
MDD & 模块开发卷宗 & 第11周 & 第12周 & 2周 & [姓名] \\
PSR & 项目开发总结报告 & 第11周 & 第12周 & 2周 & [姓名] \\
\hline
\end{longtable}

\subsubsection{里程碑计划}
\begin{longtable}{|p{3cm}|p{4cm}|p{4cm}|p{3cm}|}
\hline
里程碑 & 时间节点 & 完成标准 & 交付文档 \\
\hline
M1: 启动文档完成 & 第2周末 & FSR和PDP通过评审 & FSR v1.0, PDP v1.0 \\
M2: 需求文档完成 & 第4周末 & SRS和DRS通过评审 & SRS v1.0, DRS v1.0 \\
M3: 设计文档完成 & 第7周末 & HLD和DLD通过评审 & HLD v1.0, DLD v1.0 \\
M4: 测试文档完成 & 第9周末 & TP和TAR通过评审 & TP v1.0, TAR v1.0 \\
M5: 用户文档完成 & 第10周末 & UM和OM通过评审 & UM v1.0, OM v1.0 \\
M6: 全部文档完成 & 第12周末 & 所有文档通过评审 & 全部文档v1.0 \\
\hline
\end{longtable}

\subsubsection{进度跟踪机制}
\textbf{4.4.1 跟踪方式}
\begin{itemize}
\item 每周召开文档编制进度会议
\item 每周提交文档编制进度报告
\item 使用项目管理工具跟踪进度
\end{itemize}

\textbf{4.4.2 进度报告}
\begin{itemize}
\item 报告周期：每周
\item 报告内容：本周完成情况、下周计划、存在问题、风险识别
\item 报告格式：按照项目进度报告模板
\end{itemize}

\textbf{4.4.3 进度调整}
\begin{itemize}
\item 当进度出现偏差时，及时调整计划
\item 进度调整需经过文档负责人批准
\item 更新文档计划版本
\end{itemize}

\hrule
\textbf{注意}：文档格式和标准要求已移至独立的格式规划文档：\texttt{docs/gb8567-docs/PROJECT\_FORMAT\_MASTER\_PLAN.md}

请参考该文档了解所有格式要求、转换方法和验证标准。

\hrule
\subsection{文档评审与批准}
\subsubsection{评审流程}
\begin{itemize}
\item 优势：版本控制友好、轻量级、支持代码高亮
\item 适用：技术文档、API文档、开发文档、README

\end{itemize}
\begin{itemize}
\item \textbf{PDF格式}：约20\%的正式文档使用PDF格式
\begin{itemize}
\item 优势：格式固定、打印友好、正式性强
\item 适用：正式交付文档、归档文档、合同文档

\end{itemize}
\item \textbf{DOCX格式}：约5\%的文档使用DOCX格式
\begin{itemize}
\item 优势：编辑方便、格式丰富、适合复杂排版
\item 适用：需要复杂格式的文档、与Microsoft Office集成的场景
\end{itemize}
\end{itemize}

\textbf{5.1.3 本项目格式选择}

\textbf{主文档格式}：Markdown (.md)
\begin{itemize}
\item 理由：便于版本控制、团队协作、技术文档编写
\item 工具：VS Code、Mark Text、Typora等编辑器
\item 版本控制：Git仓库管理
\item \textbf{重要说明}：Markdown本身不支持字号、行距等格式设置，格式控制通过转换工具实现
\end{itemize}

\textbf{评审文档格式}：DOCX (.docx) [警告] \textbf{关键格式}
\begin{itemize}
\item 理由：评审环节需要边评审边修改，DOCX格式最适合
\item 转换方式：使用Pandoc将Markdown转换为DOCX（使用参考模板）
\item 参考模板：必须创建符合GB/T 8567-2006标准的参考DOCX模板（\texttt{templates/reference.docx}）
\item 转换工具：使用\texttt{tools/convert\_md\_to\_docx.bat}（Windows）或\texttt{tools/convert\_md\_to\_docx.sh}（Linux/macOS）
\item \textbf{参考文档}：详见\texttt{docs/gb8567-docs/MD\_TO\_DOCX\_SOLUTION.md}和\texttt{docs/gb8567-docs/DOCX\_TEMPLATE\_GUIDE.md}
\end{itemize}

\textbf{发布文档格式}：PDF (.pdf)
\begin{itemize}
\item 理由：正式交付、格式固定、便于归档
\item 生成方式：通过Word的"另存为PDF"功能从DOCX生成，确保格式一致性
\item \textbf{不推荐}：直接从Markdown转换PDF（格式控制不如DOCX→PDF）
\item 存储位置：{\CJKfamily{zhfs}\ttfamily docs/gb8567-docs/[分类]/pdf/}目录
\end{itemize}

\textbf{图表格式}：
\begin{itemize}
\item 流程图、架构图：SVG格式（矢量图，可缩放）
\item 截图、照片：PNG格式（无损压缩）
\item 复杂图表：PDF格式（如Excel导出的图表）
\end{itemize}

\textbf{5.1.4 格式实现方式说明}

\textbf{关键理解}：
\begin{itemize}
\item [警告] \textbf{Markdown是内容标记语言，不是格式控制语言}
\item [正确] \textbf{格式控制需要在转换阶段通过工具实现}
\item [正确] \textbf{格式规范主要针对最终发布的PDF/DOCX文档}
\end{itemize}

\textbf{格式实现流程}：

\begin{verbatim}
Markdown编写（关注内容结构）
    ↓
添加YAML元数据（格式参数）
    ↓
使用转换工具（pandoc等）
    ↓
应用格式模板/样式
    ↓
PDF/DOCX输出（符合标准格式）
    ↓
格式验证（确保符合要求）
\end{verbatim}

\textbf{格式控制方式}：

\begin{enumerate}
\item \textbf{YAML元数据方式}（推荐）：
\begin{itemize}
\item 在Markdown文件头部添加YAML元数据区块
\item 定义字体、字号、行距、页边距等参数
\item 使用pandoc转换时自动应用这些参数

\end{itemize}
\item \textbf{CSS样式方式}：
\begin{itemize}
\item 创建CSS样式文件
\item 定义所有格式要求
\item 转换时应用CSS样式

\end{itemize}
\item \textbf{模板文件方式}：
\begin{itemize}
\item 创建pandoc模板文件（.tex或.docx）
\item 定义所有格式要求
\item 转换时使用模板文件
\end{itemize}
\end{enumerate}

\textbf{格式验证}：
\begin{itemize}
\item 转换后必须验证PDF/DOCX格式是否符合GB/T 8567-2006标准
\item 使用格式检查清单进行验证
\item 必要时调整转换配置
\end{itemize}

\textbf{参考文档}：
\begin{itemize}
\item \texttt{docs/gb8567-docs/MARKDOWN\_FORMAT\_SOLUTION.md} - 格式解决方案详细说明
\item \texttt{docs/templates/MARKDOWN\_TEMPLATE\_WITH\_YAML.md} - 带YAML元数据的Markdown模板
\item \texttt{tools/convert\_md\_to\_pdf.sh} / \texttt{tools/convert\_md\_to\_pdf.bat} - 转换脚本
\end{itemize}

\subsubsection{文档结构要求（详细规范）}
\textbf{5.2.1 封面页格式}

封面页必须包含以下信息，格式如下：

\begin{verbatim}
┌─────────────────────────────────────────┐
│                                          │
│         [项目名称]                       │
│                                          │
│      [文档编号]                           │
│      [文档名称]                           │
│                                          │
│                                          │
│  编制单位：[单位名称]                     │
│  编制日期：YYYY-MM-DD                    │
│  版本号：v1.0                            │
│                                          │
│                                          │
│    [编制单位Logo（可选）]                 │
│                                          │
└─────────────────────────────────────────┘
\end{verbatim}

\textbf{格式要求}：
\begin{itemize}
\item 文档名称：黑体，小一号（24pt），居中
\item 文档编号：黑体，三号（16pt），居中
\item 项目名称：黑体，二号（22pt），居中
\item 其他信息：宋体，四号（14pt），居中
\item 页面：A4，上下左右边距均为2.5cm
\end{itemize}

\textbf{5.2.2 文档变更记录表格式}

\begin{longtable}{|p{4cm}|p{5cm}|p{5cm}|}
\hline
项目 & 格式要求 & 示例 \\
\hline
表格宽度 & 占页面宽度的90\% & - \\
表格边框 & 外边框粗线（1.5pt），内边框细线（0.5pt） & - \\
表头背景 & 浅灰色（RGB: 240,240,240） & - \\
表头字体 & 黑体，五号（10.5pt），加粗 & - \\
表格内容 & 宋体，五号（10.5pt） & - \\
行高 & 最小值0.8cm & - \\
对齐方式 & 版本号、日期居中，其他左对齐 & - \\
\hline
\end{longtable}

\textbf{5.2.3 目录格式}

\begin{itemize}
\item \textbf{生成方式}：自动生成（Markdown编辑器自动生成，或使用工具生成）
\item \textbf{格式要求}：
\begin{itemize}
\item 目录标题：黑体，三号（16pt），居中
\item 目录内容：宋体，小四号（12pt）
\item 章节编号：使用阿拉伯数字，如1、1.1、1.1.1
\item 页码：右对齐，使用点线连接
\item 行距：单倍行距
\end{itemize}
\item \textbf{示例}：
\end{itemize}
  \begin{verbatim}
  目录
\begin{enumerate}
\item 引言 ................................ 1
\end{enumerate}
     1.1 目的 ............................ 1
     1.2 范围 ............................ 2
\begin{enumerate}
\item 文档编制范围 ........................ 3
\end{enumerate}
  \end{verbatim}

\textbf{5.2.4 正文结构}

\begin{itemize}
\item \textbf{章节标题格式}：
\begin{itemize}
\item 一级标题（章节）：黑体，小二号（18pt），加粗，段前12pt，段后6pt
\item 二级标题（小节）：黑体，三号（16pt），加粗，段前6pt，段后3pt
\item 三级标题（子节）：黑体，四号（14pt），加粗，段前3pt，段后3pt
\item 四级标题（条）：黑体，小四号（12pt），加粗，段前3pt，段后0pt

\end{itemize}
\item \textbf{正文段落格式}：
\begin{itemize}
\item 字体：宋体，小四号（12pt）
\item 行距：1.5倍行距
\item 段前：0.5行（6pt）
\item 段后：0.5行（6pt）
\item 首行缩进：2字符（中文段落）
\item 对齐方式：两端对齐

\end{itemize}
\item \textbf{列表格式}：
\begin{itemize}
\item 无序列表：使用项目符号（•、○、-），悬挂缩进2字符
\item 有序列表：使用阿拉伯数字（1、2、3），悬挂缩进2字符
\item 列表项行距：单倍行距
\item 列表项间距：0.3行
\end{itemize}
\end{itemize}

\subsubsection{字体和字号详细规范}
\textbf{5.3.1 字体选择原则}

\begin{longtable}{|p{4cm}|p{5cm}|p{5cm}|}
\hline
文字类型 & 字体 & 理由 \\
\hline
\textbf{中文正文} & 宋体 & 标准中文文档字体，易读性好 \\
\textbf{中文标题} & 黑体 & 突出重点，层次分明 \\
\textbf{英文正文} & Times New Roman & 标准英文文档字体 \\
\textbf{英文标题} & Arial & 简洁现代 \\
\textbf{代码文本} & Courier New / Consolas & 等宽字体，便于对齐 \\
\textbf{数学公式} & Times New Roman & 数学公式标准字体 \\
\textbf{表格文字} & 宋体（中文）/ Times New Roman（英文） & 表格空间有限，使用标准字体 \\
\hline
\end{longtable}

\textbf{5.3.2 字号详细规范}

\begin{longtable}{|p{3cm}|p{4cm}|p{4cm}|p{3cm}|}
\hline
文字类型 & 字号 & 磅值 & 应用场景 \\
\hline
\textbf{文档标题（封面）} & 小一号 & 24pt & 封面文档名称 \\
\textbf{一级标题（章节）} & 小二号 & 18pt & 1、2、3... \\
\textbf{二级标题（小节）} & 三号 & 16pt & 1.1、1.2、1.3... \\
\textbf{三级标题（子节）} & 四号 & 14pt & 1.1.1、1.1.2... \\
\textbf{四级标题（条）} & 小四号 & 12pt & 1.1.1.1、1.1.1.2... \\
\textbf{正文} & 小四号 & 12pt & 正文段落 \\
\textbf{表格文字} & 五号 & 10.5pt & 表格内容 \\
\textbf{表格标题} & 小五号 & 9pt & 表格标题 \\
\textbf{图表标题} & 小四号 & 12pt & 图标题、表标题 \\
\textbf{图表说明} & 五号 & 10.5pt & 图表下方的说明文字 \\
\textbf{代码文本} & 五号 & 10.5pt & 代码块、行内代码 \\
\textbf{页眉页脚} & 小五号 & 9pt & 页眉、页脚文字 \\
\textbf{注释} & 小五号 & 9pt & 脚注、注释 \\
\hline
\end{longtable}

\textbf{5.3.3 字体加粗规则}

\begin{itemize}
\item \textbf{必须加粗}：所有标题、表格表头、重要术语首次出现
\item \textbf{可选加粗}：强调内容、关键词、重要数字
\item \textbf{禁止加粗}：正文段落、正文中的普通文字
\end{itemize}

\subsubsection{段落格式详细规范}
\textbf{5.4.1 行距规范}

\begin{longtable}{|p{4cm}|p{5cm}|p{5cm}|}
\hline
内容类型 & 行距 & 说明 \\
\hline
\textbf{正文段落} & 1.5倍行距 & 标准行距，易读性好 \\
\textbf{标题与正文之间} & 单倍行距 & 标题段后+正文段前=适当间距 \\
\textbf{列表项} & 单倍行距 & 列表项之间紧凑 \\
\textbf{表格内容} & 单倍行距 & 表格空间有限 \\
\textbf{代码块} & 单倍行距 & 代码通常行数较多 \\
\hline
\end{longtable}

\textbf{5.4.2 段间距规范}

\begin{longtable}{|p{3cm}|p{4cm}|p{4cm}|p{3cm}|}
\hline
段落类型 & 段前 & 段后 & 说明 \\
\hline
\textbf{一级标题} & 12pt & 6pt & 与上一章节分隔明显 \\
\textbf{二级标题} & 6pt & 3pt & 与上一小节分隔 \\
\textbf{三级标题} & 3pt & 3pt & 适当的间距 \\
\textbf{正文段落} & 0.5行（6pt） & 0.5行（6pt） & 标准段落间距 \\
\textbf{列表项} & 0.3行（3.6pt） & 0.3行（3.6pt） & 列表项之间紧凑 \\
\textbf{表格前后} & 6pt & 6pt & 表格与正文的间距 \\
\hline
\end{longtable}

\textbf{5.4.3 缩进规范}

\begin{itemize}
\item \textbf{首行缩进}：
\begin{itemize}
\item 中文正文段落：2字符（7.5mm）
\item 英文正文段落：0（不使用首行缩进）
\item 列表项：悬挂缩进2字符

\end{itemize}
\item \textbf{段落缩进}：
\begin{itemize}
\item 正文段落：左右缩进均为0
\item 引用段落：左右各缩进2字符，使用特殊格式（如浅灰色背景）
\end{itemize}
\end{itemize}

\subsubsection{页面设置详细规范}
\textbf{5.5.1 页面规格}

\begin{itemize}
\item \textbf{纸张大小}：A4（210mm × 297mm）
\item \textbf{方向}：纵向（Portrait）
\item \textbf{页边距}：
\begin{itemize}
\item 上边距：2.5cm（用于页眉）
\item 下边距：2.5cm（用于页脚）
\item 左边距：2.5cm（装订线）
\item 右边距：2.5cm
\end{itemize}
\end{itemize}

\textbf{5.5.2 页眉页脚格式}

\textbf{页眉格式}：
\begin{itemize}
\item 位置：页眉高度1.5cm，距离页面顶部1.25cm
\item 左侧：文档名称（黑体，小五号，9pt）
\item 右侧：版本号（黑体，小五号，9pt）
\item 分隔线：页眉下方0.5pt细线，颜色RGB(200,200,200)
\end{itemize}

\textbf{页脚格式}：
\begin{itemize}
\item 位置：页脚高度1.5cm，距离页面底部1.25cm
\item 左侧：编制单位名称（宋体，小五号，9pt）
\item 中间：页码（宋体，小五号，9pt），格式：第X页 共Y页
\item 右侧：编制日期（宋体，小五号，9pt），格式：YYYY-MM-DD
\end{itemize}

\textbf{5.5.3 页码格式}

\begin{itemize}
\item \textbf{编号方式}：阿拉伯数字（1、2、3...）
\item \textbf{起始页码}：目录页从罗马数字开始（i、ii、iii...），正文从1开始
\item \textbf{位置}：页脚居中
\item \textbf{格式}：第X页 共Y页（如：第1页 共50页）
\end{itemize}

\subsubsection{编号规则详细规范}
\textbf{5.6.1 章节编号规则}

\begin{itemize}
\item \textbf{编号方式}：阿拉伯数字分级编号
\item \textbf{编号格式}：
\begin{itemize}
\item 一级：1、2、3...
\item 二级：1.1、1.2、1.3...
\item 三级：1.1.1、1.1.2、1.1.3...
\item 四级：1.1.1.1、1.1.1.2...（仅在必要时使用）

\end{itemize}
\item \textbf{编号规则}：
\begin{itemize}
\item 章节编号后必须有一个空格，再跟标题文字
\item 示例：{\CJKfamily{zhfs}\ttfamily 1. 引言}、{\CJKfamily{zhfs}\ttfamily 1.1 目的}、{\CJKfamily{zhfs}\ttfamily 1.1.1 文档目的}
\end{itemize}
\end{itemize}

\textbf{5.6.2 图表编号规则}

\textbf{图表编号格式}：
\begin{itemize}
\item 图形：{\CJKfamily{zhfs}\ttfamily 图[章节号]-[序号]}，如{\CJKfamily{zhfs}\ttfamily 图1-1}、{\CJKfamily{zhfs}\ttfamily 图2-3}
\item 表格：{\CJKfamily{zhfs}\ttfamily 表[章节号]-[序号]}，如{\CJKfamily{zhfs}\ttfamily 表1-1}、{\CJKfamily{zhfs}\ttfamily 表2-3}
\end{itemize}

\textbf{图表标题格式}：
\begin{itemize}
\item 图形标题：位于图形下方，居中，格式：{\CJKfamily{zhfs}\ttfamily 图X-Y [图形标题]}
\item 表格标题：位于表格上方，居中，格式：{\CJKfamily{zhfs}\ttfamily 表X-Y [表格标题]}
\item 字体：黑体，小四号（12pt）
\end{itemize}

\textbf{图表说明格式}：
\begin{itemize}
\item 位置：图表标题下方
\item 格式：宋体，五号（10.5pt），左对齐
\item 内容：简要说明图表的内容、数据来源、用途等
\end{itemize}

\textbf{5.6.3 公式编号规则}

\begin{itemize}
\item \textbf{编号格式}：{\CJKfamily{zhfs}\ttfamily ([章节号]-[序号])}，如\texttt{(1-1)}、\texttt{(2-3)}
\item \textbf{位置}：公式右侧，右对齐
\item \textbf{字体}：Times New Roman，五号（10.5pt）
\item \textbf{示例}：
\end{itemize}
  \begin{lstlisting}
  E = mc²                                    (1-1)
  \end{lstlisting}

\textbf{5.6.4 引用格式规则}

\textbf{引用标准}：
\begin{itemize}
\item 格式：{\CJKfamily{zhfs}\ttfamily GB/T 8567-2006《计算机软件文档编制规范》}
\item 首次出现：使用全称
\item 再次出现：可使用简称\texttt{GB/T 8567-2006}
\end{itemize}

\textbf{引用文档}：
\begin{itemize}
\item 格式：{\CJKfamily{zhfs}\ttfamily [文档编号]}或{\CJKfamily{zhfs}\ttfamily （文档编号）}
\item 示例：\texttt{[FSR-001]}、{\CJKfamily{zhfs}\ttfamily （可行性研究报告）}
\end{itemize}

\textbf{引用图表}：
\begin{itemize}
\item 格式：{\CJKfamily{zhfs}\ttfamily 见图X-Y}、{\CJKfamily{zhfs}\ttfamily 如表X-Y所示}
\item 示例：{\CJKfamily{zhfs}\ttfamily 见图1-1}、{\CJKfamily{zhfs}\ttfamily 如表2-3所示}
\end{itemize}

\subsubsection{表格格式详细规范}
\textbf{5.7.1 表格基本格式}

\begin{longtable}{|p{4cm}|p{5cm}|p{5cm}|}
\hline
项目 & 格式要求 & 说明 \\
\hline
\textbf{表格宽度} & 占页面宽度的90\%-100\% & 根据内容调整 \\
\textbf{表格边框} & 外边框1.5pt粗线，内边框0.5pt细线 & 使用实线 \\
\textbf{表格对齐} & 居中 & 表格在页面中居中 \\
\textbf{表头背景} & RGB(240,240,240)浅灰色 & 突出表头 \\
\textbf{表头字体} & 黑体，五号（10.5pt），加粗 & 突出表头 \\
\textbf{表格内容字体} & 宋体（中文）/ Times New Roman（英文），五号（10.5pt） & 标准字体 \\
\textbf{行高} & 最小值0.8cm，根据内容调整 & 保证可读性 \\
\textbf{单元格内边距} & 上下0.2cm，左右0.3cm & 适当间距 \\
\hline
\end{longtable}

\textbf{5.7.2 表格内容对齐}

\begin{longtable}{|p{4cm}|p{5cm}|p{5cm}|}
\hline
内容类型 & 对齐方式 & 说明 \\
\hline
\textbf{文字内容} & 左对齐 & 标准对齐方式 \\
\textbf{数字内容} & 右对齐 & 便于比较数值 \\
\textbf{序号} & 居中 & 居中美观 \\
\textbf{日期} & 居中 & 居中美观 \\
\textbf{表头} & 居中 & 表头居中 \\
\hline
\end{longtable}

\textbf{5.7.3 表格跨页处理}

\begin{itemize}
\item \textbf{表头重复}：跨页表格每页重复表头
\item \textbf{续表标识}：第二页开始，表格上方标注"续表X-Y"
\item \textbf{格式}：宋体，小五号（9pt），居中
\end{itemize}

\subsubsection{图表格式详细规范}
\textbf{5.8.1 图形格式要求}

\begin{longtable}{|p{4cm}|p{5cm}|p{5cm}|}
\hline
项目 & 格式要求 & 说明 \\
\hline
\textbf{图形格式} & SVG（优先）/ PNG（次选） & SVG矢量图可缩放 \\
\textbf{图形分辨率} & 至少300dpi（打印质量） & 确保清晰度 \\
\textbf{图形大小} & 宽度不超过页面宽度的90\% & 适应页面 \\
\textbf{图形颜色} & 打印友好（避免纯黄色等） & 确保打印清晰 \\
\textbf{图形边框} & 可选，0.5pt细线 & 根据需要添加 \\
\hline
\end{longtable}

\textbf{5.8.2 图形标题和说明}

\begin{itemize}
\item \textbf{标题位置}：图形下方
\item \textbf{标题格式}：{\CJKfamily{zhfs}\ttfamily 图X-Y [图形标题]}，黑体，小四号（12pt），居中
\item \textbf{说明位置}：标题下方
\item \textbf{说明格式}：宋体，五号（10.5pt），左对齐，首行缩进2字符
\end{itemize}

\textbf{5.8.3 图形引用}

\begin{itemize}
\item \textbf{引用格式}：{\CJKfamily{zhfs}\ttfamily 见图X-Y}、{\CJKfamily{zhfs}\ttfamily 如图X-Y所示}、{\CJKfamily{zhfs}\ttfamily 参见图X-Y}
\item \textbf{引用位置}：正文中引用图形时，使用上述格式
\item \textbf{示例}：{\CJKfamily{zhfs}\ttfamily 系统架构见图2-1所示。}
\end{itemize}

\subsubsection{代码格式详细规范}
\textbf{5.9.1 代码块格式}

\begin{itemize}
\item \textbf{字体}：Courier New或Consolas（等宽字体）
\item \textbf{字号}：五号（10.5pt）
\item \textbf{背景色}：浅灰色（RGB(245,245,245)）
\item \textbf{边框}：1pt细线，颜色RGB(200,200,200)
\item \textbf{内边距}：上下0.3cm，左右0.5cm
\item \textbf{行距}：单倍行距
\item \textbf{缩进}：使用空格，不使用Tab（设置为4个空格）
\end{itemize}

\textbf{5.9.2 行内代码格式}

\begin{itemize}
\item \textbf{格式}：使用反引号包裹，如\texttt{code}
\item \textbf{字体}：Courier New，五号（10.5pt）
\item \textbf{背景色}：浅灰色（RGB(245,245,245)）
\item \textbf{边框}：无
\end{itemize}

\textbf{5.9.3 代码语言标识}

\begin{itemize}
\item \textbf{格式}：代码块开头使用语言标识
\item \textbf{示例}：
\end{itemize}
  \begin{lstlisting}[language=python]
  def hello():
      print("Hello, World!")
  \end{lstlisting}

\subsubsection{文档命名规范（详细）}
\textbf{5.10.1 文件命名格式}

\begin{verbatim}
[文档编号]\_[文档名称]\_v[版本号].[扩展名]
\end{verbatim}

\textbf{详细规则}：
\begin{itemize}
\item \textbf{文档编号}：按照GB/T 8567-2006标准（如FSR、PDP、SRS）
\item \textbf{文档名称}：中文全称，与标准文档名称一致
\item \textbf{版本号}：格式{\CJKfamily{zhfs}\ttfamily v[主版本号].[次版本号].[修订号]}
\begin{itemize}
\item 主版本号：重大变更（如v1.0 → v2.0）
\item 次版本号：功能增加（如v1.0 → v1.1）
\item 修订号：错误修正（如v1.0 → v1.0.1）
\end{itemize}
\item \textbf{扩展名}：
\begin{itemize}
\item \texttt{.md}：Markdown源文件
\item \texttt{.pdf}：PDF发布文件
\item \texttt{.docx}：Word文档（如需）
\end{itemize}
\end{itemize}

\textbf{示例}：
\begin{itemize}
\item {\CJKfamily{zhfs}\ttfamily FSR\_可行性研究报告\_v1.0.md}
\item {\CJKfamily{zhfs}\ttfamily PDP\_项目开发计划\_v1.1.md}
\item {\CJKfamily{zhfs}\ttfamily SRS\_软件需求说明书\_v1.0.1.md}
\end{itemize}

\textbf{5.10.2 文件命名规则说明}

\begin{itemize}
\item \textbf{避免使用特殊字符}：文件名中不使用\texttt{/ \ : * ? " < > |}等特殊字符
\item \textbf{使用下划线分隔}：使用下划线\texttt{\_}分隔各部分
\item \textbf{保持一致性}：所有文档使用相同的命名规则
\item \textbf{版本号格式}：版本号使用小写\texttt{v}，如\texttt{v1.0}而非\texttt{V1.0}
\end{itemize}

\subsubsection{文档编号规则（详细）}
\textbf{5.11.1 文档编号格式}

\begin{verbatim}
[项目编号]-[文档类型编号]-[序号]
\end{verbatim}

\textbf{详细规则}：
\begin{itemize}
\item \textbf{项目编号}：项目唯一标识符（如PROJ、JMETER-TEST）
\item \textbf{文档类型编号}：按照GB/T 8567-2006标准
\begin{itemize}
\item FSR：可行性研究报告
\item PDP：项目开发计划
\item SRS：软件需求说明书
\item DRS：数据要求说明书
\item HLD：概要设计说明书
\item DLD：详细设计说明书
\item DBD：数据库设计说明书
\item MDD：模块开发卷宗
\item TP：测试计划
\item TAR：测试分析报告
\item UM：用户手册
\item OM：操作手册
\item MPR：开发进度月报
\item PSR：项目开发总结报告
\end{itemize}
\item \textbf{序号}：同一类型文档的序号，使用三位数字（001、002...）
\end{itemize}

\textbf{示例}：
\begin{itemize}
\item \texttt{PROJ-FSR-001}：项目可行性研究报告
\item \texttt{PROJ-PDP-001}：项目开发计划
\item \texttt{PROJ-SRS-001}：软件需求说明书
\end{itemize}

\textbf{5.11.2 文档编号使用规则}

\begin{itemize}
\item \textbf{文档编号出现在}：
\begin{itemize}
\item 封面页
\item 页眉（可选）
\item 文档变更记录表
\item 文档索引

\end{itemize}
\item \textbf{文档编号格式}：
\begin{itemize}
\item 在封面页：使用较大字号（四号，14pt）
\item 在其他位置：使用标准字号（五号，10.5pt）
\end{itemize}
\end{itemize}

\subsubsection{文档样式统一规范}
\textbf{5.12.1 术语使用规范}

\begin{itemize}
\item \textbf{术语一致性}：
\begin{itemize}
\item 同一术语在全文档中保持一致
\item 首次出现时给出英文对照和定义
\item 建立项目术语表，统一管理

\end{itemize}
\item \textbf{术语格式}：
\begin{itemize}
\item 首次出现：{\CJKfamily{zhfs}\ttfamily 术语（英文）}，如{\CJKfamily{zhfs}\ttfamily 可行性研究报告（Feasibility Study Report）}
\item 再次出现：直接使用中文术语
\item 英文术语：首字母大写，如\texttt{Feasibility Study Report}
\end{itemize}
\end{itemize}

\textbf{5.12.2 标点符号规范}

\begin{itemize}
\item \textbf{中文标点}：中文文档使用中文标点符号
\begin{itemize}
\item 句号：\texttt{。}
\item 逗号：\texttt{，}
\item 分号：\texttt{；}
\item 冒号：\texttt{：}
\item 引号：\texttt{""}（中文引号）

\end{itemize}
\item \textbf{英文标点}：英文内容使用英文标点符号
\begin{itemize}
\item 句号：\texttt{.}
\item 逗号：\texttt{,}
\item 分号：\texttt{;}
\item 冒号：\texttt{:}
\item 引号：\texttt{""}（英文引号）

\end{itemize}
\item \textbf{数字和标点}：
\begin{itemize}
\item 数字与中文之间：使用中文标点，如{\CJKfamily{zhfs}\ttfamily 项目共14个文档。}
\item 数字与英文之间：使用英文标点，如\texttt{The project has 14 documents.}
\end{itemize}
\end{itemize}

\textbf{5.12.3 数字使用规范}

\begin{itemize}
\item \textbf{阿拉伯数字}：
\begin{itemize}
\item 用于：版本号、日期、页码、图表编号、公式编号、技术参数
\item 示例：\texttt{v1.0}、\texttt{2025-11-05}、{\CJKfamily{zhfs}\ttfamily 第1页}、{\CJKfamily{zhfs}\ttfamily 图1-1}、\texttt{(1-1)}

\end{itemize}
\item \textbf{中文数字}：
\begin{itemize}
\item 用于：章节编号、一般性描述、习惯用法
\item 示例：{\CJKfamily{zhfs}\ttfamily 第一章}、{\CJKfamily{zhfs}\ttfamily 十二周}、{\CJKfamily{zhfs}\ttfamily 三个月}

\end{itemize}
\item \textbf{数字格式}：
\begin{itemize}
\item 日期：\texttt{YYYY-MM-DD}格式，如\texttt{2025-11-05}
\item 版本号：{\CJKfamily{zhfs}\ttfamily v主版本号.次版本号.修订号}，如\texttt{v1.0}、\texttt{v1.1}、\texttt{v1.0.1}
\item 百分比：使用百分号，如\texttt{75\%}
\end{itemize}
\end{itemize}

\textbf{5.12.4 单位使用规范}

\begin{itemize}
\item \textbf{单位格式}：
\begin{itemize}
\item 中文单位：使用中文单位，如{\CJKfamily{zhfs}\ttfamily 厘米}、{\CJKfamily{zhfs}\ttfamily 毫米}、{\CJKfamily{zhfs}\ttfamily 磅}
\item 英文单位：使用英文单位，如\texttt{cm}、\texttt{mm}、\texttt{pt}
\item 混合使用：数字与单位之间空一格，如\texttt{12 pt}、\texttt{2.5 cm}

\end{itemize}
\item \textbf{常用单位}：
\begin{itemize}
\item 长度：\texttt{cm}（厘米）、\texttt{mm}（毫米）、\texttt{pt}（磅）
\item 字号：\texttt{pt}（磅）或{\CJKfamily{zhfs}\ttfamily 号}（字号）
\item 颜色：RGB格式，如\texttt{RGB(240,240,240)}
\end{itemize}
\end{itemize}

\subsubsection{Markdown转DOCX格式规范（评审文档）}
\textbf{5.13.1 评审文档格式要求}

\textbf{重要说明}：
\begin{itemize}
\item [警告] \textbf{评审环节必须使用DOCX格式}：真实项目中的评审环节需要边评审边修改，DOCX格式最适合
\item [警告] \textbf{PDF由DOCX导出}：PDF作为最终交付物，通过Word的"另存为PDF"功能生成，确保格式一致性
\item [警告] \textbf{避免维护两套文档}：只维护Markdown源文件，通过工具自动转换为DOCX
\end{itemize}

\textbf{工作流程}：
\begin{verbatim}
Markdown编写（源文件）
    ↓
Pandoc转换（使用参考DOCX模板）
    ↓
DOCX文档（用于评审和修改）
    ↓
Word导出PDF（最终交付物）
\end{verbatim}

\textbf{转换工具}：
\begin{itemize}
\item \textbf{Windows}：\texttt{tools/convert\_md\_to\_docx.bat}
\item \textbf{Linux/macOS}：\texttt{tools/convert\_md\_to\_docx.sh}
\item \textbf{直接命令}：\texttt{pandoc input.md -o output.docx --reference-doc=templates/reference.docx}
\end{itemize}

\textbf{参考文档}：
\begin{itemize}
\item \texttt{docs/gb8567-docs/MD\_TO\_DOCX\_SOLUTION.md} - 完整解决方案
\item \texttt{docs/gb8567-docs/DOCX\_TEMPLATE\_GUIDE.md} - 参考模板创建指南
\end{itemize}

\textbf{5.13.2 参考DOCX模板要求}

参考DOCX模板必须包含以下样式定义：

\begin{enumerate}
\item \textbf{正文样式（Normal）}：
\begin{itemize}
\item 字体：宋体（中文）、Times New Roman（英文）
\item 字号：小四号（12pt）
\item 行距：1.5倍
\item 段前距：0pt
\item 段后距：6pt

\end{itemize}
\item \textbf{标题样式}：
\begin{itemize}
\item 一级标题（Heading 1）：黑体、三号（16pt）、加粗
\item 二级标题（Heading 2）：黑体、四号（14pt）、加粗
\item 三级标题（Heading 3）：黑体、小四号（12pt）、加粗

\end{itemize}
\item \textbf{表格样式}：
\begin{itemize}
\item 表格边框：实线、0.5pt
\item 表头：加粗、居中
\item 表格字体：宋体、五号（10.5pt）

\end{itemize}
\item \textbf{页面设置}：
\begin{itemize}
\item 纸张：A4
\item 页边距：上下2.5cm，左右2.5cm

\end{itemize}
\item \textbf{页眉页脚}：
\begin{itemize}
\item 页眉：文档名称、版本号
\item 页脚：页码（居中）
\end{itemize}
\end{enumerate}

\textbf{详细创建步骤}：请参考\texttt{docs/gb8567-docs/DOCX\_TEMPLATE\_GUIDE.md}

\textbf{5.13.3 转换验证清单}

转换完成后，请检查：

\begin{itemize}
\item [ ] 字体和字号是否正确（正文：宋体12pt）
\item [ ] 行距是否正确（1.5倍行距）
\item [ ] 页边距是否正确（2.5cm）
\item [ ] 标题格式是否正确（黑体、加粗）
\item [ ] 表格格式是否正确（边框、对齐）
\item [ ] 页眉页脚是否正确
\item [ ] 目录是否正确生成
\item [ ] 可以在Word中正常编辑
\item [ ] 可以添加批注
\item [ ] 可以导出为PDF
\end{itemize}

\subsubsection{Markdown转PDF格式规范（最终交付）}
\textbf{5.14.1 转换工具选择}

\textbf{重要说明}：
\begin{itemize}
\item [警告] \textbf{PDF由DOCX导出}：不直接从Markdown转换PDF，而是通过Word的"另存为PDF"功能从DOCX生成
\item [正确] \textbf{格式一致性}：确保PDF格式与DOCX格式完全一致
\item [正确] \textbf{评审流程}：DOCX用于评审，PDF用于最终交付
\end{itemize}

\textbf{PDF生成方式}：
\begin{enumerate}
\item 在Word中打开转换后的DOCX文档
\item 检查格式是否符合要求
\item 点击 \textbf{文件} → \textbf{另存为} → \textbf{PDF}
\item 选择PDF选项（保持格式、可搜索等）
\item 保存PDF文件
\end{enumerate}

\textbf{5.14.2 转换工具选择（备用方案）}

如果需要直接从Markdown转换PDF（不推荐，仅作为备用方案），可以使用以下工具：

\begin{longtable}{|p{4cm}|p{5cm}|p{5cm}|}
\hline
工具 & 优势 & 适用场景 \\
\hline
\textbf{pandoc} & 功能强大，支持多种格式 & 批量转换、自动化流程 \\
\textbf{Mark Text} & 界面友好，实时预览 & 手动转换、单文档转换 \\
\textbf{Typora} & 所见即所得，格式丰富 & 复杂格式文档 \\
\textbf{VS Code插件} & 集成编辑器中 & 开发环境集成 \\
\hline
\end{longtable}

\textbf{5.14.3 转换配置要求}（备用方案）

\textbf{重要说明}：
\begin{itemize}
\item [警告] 以下格式要求\textbf{需要在转换时通过工具实现}
\item [警告] Markdown源文件本身\textbf{无法存储}这些格式信息
\item [正确] 必须在转换前\textbf{准备好转换配置}（YAML元数据、CSS样式、模板文件）
\end{itemize}

\textbf{配置方式1：YAML元数据}（推荐）

在Markdown文件头部添加YAML元数据：

\begin{verbatim}
\hrule
pdf-engine: xelatex
mainfont: "SimSun"          # 宋体（中文正文）
sansfont: "SimHei"          # 黑体（中文标题）
fontsize: 12pt              # 正文字号：小四号（12pt）
linestretch: 1.5            # 行距：1.5倍
geometry: margin=2.5cm       # 页边距：2.5cm
\hrule
\end{verbatim}

\textbf{配置方式2：转换命令参数}

\begin{lstlisting}[language=bash]
pandoc input.md -o output.pdf \textbackslash\{\}
  --pdf-engine=xelatex \textbackslash\{\}
  -V mainfont="SimSun" \textbackslash\{\}
  -V fontsize=12pt \textbackslash\{\}
  -V linestretch=1.5 \textbackslash\{\}
  -V geometry:margin=2.5cm
\end{lstlisting}

\textbf{配置方式3：模板文件}

使用pandoc模板文件（.tex）或参考DOCX文件（.docx）定义格式。

\textbf{格式配置详细要求}：

\begin{itemize}
\item \textbf{页面设置}：
\begin{itemize}
\item 纸张：A4
\item 页边距：上下2.5cm，左右2cm
\item 方向：纵向

\end{itemize}
\item \textbf{字体映射}：
\begin{itemize}
\item 中文正文：宋体（SimSun）
\item 中文标题：黑体（SimHei）
\item 英文正文：Times New Roman
\item 英文标题：Arial
\item 代码文本：Courier New或Consolas

\end{itemize}
\item \textbf{样式设置}：
\begin{itemize}
\item 标题样式：按照5.3节规定的字号设置
\item 正文样式：按照5.4节规定的段落格式设置
\item 表格样式：按照5.7节规定的表格格式设置
\end{itemize}
\end{itemize}

\textbf{转换工具配置}：

\begin{itemize}
\item \textbf{pandoc}：使用YAML元数据或命令行参数
\item \textbf{Mark Text}：在导出设置中配置格式
\item \textbf{Typora}：使用自定义CSS样式文件
\end{itemize}

\textbf{5.14.4 PDF输出质量要求}

\begin{itemize}
\item \textbf{分辨率}：至少300dpi（打印质量）
\item \textbf{文件大小}：优化后不超过10MB（单文档）
\item \textbf{可搜索性}：文本可选择和搜索
\item \textbf{书签}：自动生成书签（目录）
\item \textbf{超链接}：保持超链接可点击
\end{itemize}

\subsubsection{文档格式检查清单}
\textbf{5.15.1 格式检查项}

\textbf{重要说明}：
\begin{itemize}
\item [警告] 以下格式检查项\textbf{主要针对最终发布的PDF/DOCX文档}
\item [警告] Markdown源文件主要检查\textbf{内容结构}，格式检查在转换后进行
\end{itemize}

\textbf{Markdown源文件检查项}（编写阶段）：

\begin{itemize}
\item [ ] 章节编号符合规范（1、1.1、1.1.1）
\item [ ] 文档结构完整（封面、目录、正文、附录）
\item [ ] YAML元数据正确（格式参数定义）
\item [ ] 图表编号符合规范（图1-1、表2-3）
\item [ ] 文档命名符合规范
\item [ ] 文档编号符合规范
\item [ ] 术语使用一致
\item [ ] 引用格式正确
\end{itemize}

\textbf{PDF/DOCX文档检查项}（发布阶段）：

\begin{itemize}
\item [ ] 封面页格式正确（文档编号、名称、版本号等）
\item [ ] 文档变更记录表格式正确
\item [ ] 目录自动生成且格式正确
\item [ ] 字体字号符合规定（标题黑体、正文宋体）
\item [ ] 段落格式符合规定（行距1.5倍、首行缩进2字符）
\item [ ] 表格格式符合规定（边框、对齐、表头背景）
\item [ ] 页眉页脚格式正确
\item [ ] 页码格式正确（第X页 共Y页）
\item [ ] 页面设置正确（A4纸张、页边距2.5cm）
\item [ ] 打印效果验证（打印预览检查）
\end{itemize}

\textbf{5.15.2 质量检查项}

\begin{itemize}
\item [ ] 术语使用一致
\item [ ] 标点符号使用正确
\item [ ] 数字格式统一
\item [ ] 单位使用规范
\item [ ] 引用格式正确
\item [ ] 代码格式正确
\item [ ] 图表清晰美观
\end{itemize}

\subsubsection{评审流程}
文档评审分为以下阶段：
\begin{itemize}
\item 执行人：文档编写人员
\item 时间：文档编写完成后1天内
\item 内容：检查格式、内容完整性、错别字等
\end{itemize}

\textbf{6.1.2 互审}
\begin{itemize}
\item 执行人：同级文档编写人员
\item 时间：自审完成后1天内
\item 内容：交叉检查文档内容、格式等
\end{itemize}

\textbf{6.1.3 技术评审}
\begin{itemize}
\item 执行人：技术负责人
\item 时间：互审完成后2天内
\item 内容：检查技术内容准确性、完整性
\end{itemize}

\textbf{6.1.4 最终评审}
\begin{itemize}
\item 执行人：文档负责人
\item 时间：技术评审完成后1天内
\item 内容：检查文档是否符合标准要求、质量是否达标
\end{itemize}

\subsubsection{评审标准}
\begin{longtable}{|p{3cm}|p{4cm}|p{4cm}|p{3cm}|}
\hline
评审项 & 评审标准 & 权重 & 评分标准 \\
\hline
内容完整性 & 是否包含所有必需章节 & 30\% & 0-100分 \\
技术准确性 & 技术内容是否正确 & 25\% & 0-100分 \\
格式规范性 & 是否符合标准格式 & 20\% & 0-100分 \\
可读性 & 是否清晰易懂 & 15\% & 0-100分 \\
一致性 & 是否与其他文档一致 & 10\% & 0-100分 \\
\hline
\end{longtable}

\textbf{评审通过标准}：
\begin{itemize}
\item 总分≥80分
\item 各项得分≥60分
\end{itemize}

\subsubsection{批准流程}
\textbf{6.3.1 批准权限}
\begin{itemize}
\item 项目管理文档：项目负责人批准
\item 技术文档：技术负责人批准
\item 测试文档：测试负责人批准
\item 用户文档：产品负责人批准
\end{itemize}

\textbf{6.3.2 批准流程}
\begin{enumerate}
\item 文档通过最终评审后，提交批准
\item 批准人在2天内完成审批
\item 批准通过后，文档正式发布
\item 批准未通过，返回修改
\end{enumerate}

\subsubsection{评审记录}
\textbf{6.4.1 评审记录表}

\begin{longtable}{|p{2.29cm}|p{2.29cm}|p{2.29cm}|p{2.29cm}|p{2.29cm}|p{2.29cm}|p{2.29cm}|}
\hline
文档名称 & 版本号 & 评审日期 & 评审人 & 评审意见 & 评审结论 & 备注 \\
\hline
[文档名称] & v1.0 & YYYY-MM-DD & [姓名] & [意见] & 通过/不通过 & [备注] \\
\hline
\end{longtable}

\textbf{6.4.2 评审问题跟踪}
\begin{itemize}
\item 记录所有评审问题
\item 跟踪问题修改情况
\item 确认问题已解决
\end{itemize}

\hrule
\subsection{文档版本控制}
\subsubsection{版本号规则}
\textbf{7.1.1 版本号格式}
\begin{verbatim}
v[主版本号].[次版本号].[修订号]
\end{verbatim}

\textbf{示例}：
\begin{itemize}
\item v1.0：初始版本
\item v1.1：功能增加版本
\item v1.0.1：错误修正版本
\item v2.0：重大变更版本
\end{itemize}

\textbf{7.1.2 版本号规则说明}
\begin{itemize}
\item 主版本号：重大变更（架构变更、内容重构）
\item 次版本号：功能增加（新增章节、内容扩展）
\item 修订号：错误修正（错别字、格式错误）
\end{itemize}

\subsubsection{版本管理流程}
\textbf{7.2.1 版本创建}
\begin{itemize}
\item 文档初稿完成后，创建v1.0版本
\item 每次修改后，创建新版本
\item 版本号按照版本号规则递增
\end{itemize}

\textbf{7.2.2 版本存储}
\begin{itemize}
\item 使用Git进行版本管理
\item 每个版本必须提交到Git仓库
\item 建立版本标签（tag）
\end{itemize}

\textbf{7.2.3 版本发布}
\begin{itemize}
\item 文档通过评审和批准后，正式发布
\item 发布版本不可修改
\item 如需修改，创建新版本
\end{itemize}

\subsubsection{版本发布}
\textbf{7.3.1 发布流程}
\begin{enumerate}
\item 文档通过评审和批准
\item 创建发布版本标签
\item 生成发布文档（PDF格式）
\item 更新文档索引
\end{enumerate}

\textbf{7.3.2 发布记录}

\begin{longtable}{|p{2.5cm}|p{3cm}|p{3cm}|p{3cm}|p{2.5cm}|}
\hline
文档名称 & 版本号 & 发布日期 & 发布人 & 变更说明 \\
\hline
[文档名称] & v1.0 & YYYY-MM-DD & [姓名] & [变更说明] \\
\hline
\end{longtable}

\hrule
\subsection{文档存储与管理}
\subsubsection{存储结构}
文档存储结构如下：

\begin{verbatim}
docs/
├── gb8567-docs/              # GB/T 8567-2006标准文档目录
│   ├── 01-project-startup/  # 项目启动文档
│   ├── 02-requirements/     # 需求文档
│   ├── 03-design/           # 设计文档
│   ├── 04-testing/          # 测试文档
│   ├── 05-user-docs/        # 用户文档
│   └── 06-project-mgmt/     # 项目管理文档
├── templates/                # 文档模板
└── archive/                  # 文档归档
\end{verbatim}

\subsubsection{访问权限}
\begin{longtable}{|p{4cm}|p{5cm}|p{5cm}|}
\hline
角色 & 权限 & 说明 \\
\hline
文档负责人 & 读写所有文档 & 可以创建、修改、删除所有文档 \\
技术负责人 & 读写技术文档 & 可以创建、修改技术相关文档 \\
测试负责人 & 读写测试文档 & 可以创建、修改测试相关文档 \\
文档编写人员 & 读写分配文档 & 可以创建、修改分配的文档 \\
项目成员 & 只读所有文档 & 可以查看所有文档 \\
\hline
\end{longtable}

\subsubsection{备份与归档}
\textbf{8.3.1 备份策略}
\begin{itemize}
\item 每日自动备份到Git仓库
\item 每周备份到本地服务器
\item 每月备份到云端存储
\end{itemize}

\textbf{8.3.2 归档策略}
\begin{itemize}
\item 项目结束后，归档所有文档
\item 归档文档不再修改
\item 归档文档保留至少5年
\end{itemize}

\hrule
\subsection{文档质量保证}
\subsubsection{质量要求}
\textbf{9.1.1 内容质量}
\begin{itemize}
\item 完整性：包含所有必需章节
\item 准确性：技术内容准确无误
\item 一致性：术语使用一致
\item 可读性：语言清晰易懂
\end{itemize}

\textbf{9.1.2 格式质量}
\begin{itemize}
\item 符合标准格式要求
\item 格式统一规范
\item 图表清晰美观
\end{itemize}

\textbf{9.1.3 过程质量}
\begin{itemize}
\item 按照流程编制
\item 通过评审和批准
\item 版本管理规范
\end{itemize}

\subsubsection{质量检查}
\textbf{9.2.1 检查方式}
\begin{itemize}
\item 自检：文档编写人员自检
\item 互检：同级人员交叉检查
\item 专检：专职质量检查人员检查
\end{itemize}

\textbf{9.2.2 检查内容}
\begin{itemize}
\item 内容完整性检查
\item 格式规范性检查
\item 一致性检查
\item 可读性检查
\end{itemize}

\subsubsection{质量改进}
\textbf{9.3.1 问题收集}
\begin{itemize}
\item 收集评审中发现的问题
\item 收集使用中发现的问题
\item 收集用户反馈
\end{itemize}

\textbf{9.3.2 改进措施}
\begin{itemize}
\item 分析问题原因
\item 制定改进措施
\item 更新文档计划
\item 持续改进
\end{itemize}

\hrule
\subsection{附录}
\subsubsection{附录A：文档模板索引}
\begin{longtable}{|p{4cm}|p{5cm}|p{5cm}|}
\hline
文档编号 & 文档名称 & 模板文件路径 \\
\hline
FSR & 可行性研究报告 & docs/templates/FSR\_Template.md \\
PDP & 项目开发计划 & docs/templates/PDP\_Template.md \\
SRS & 软件需求说明书 & docs/templates/SRS\_Template.md \\
... & ... & ... \\
\hline
\end{longtable}

\subsubsection{附录B：参考标准}
\begin{itemize}
\item GB/T 8567-2006《计算机软件文档编制规范》
\item GB/T 9385-2008《计算机软件需求规格说明规范》
\item GB/T 9386-2008《计算机软件测试文档编制规范》
\item GB/T 16260-2006《软件工程 产品质量》
\end{itemize}

\subsubsection{附录C：术语表}
\begin{longtable}{|p{4cm}|p{5cm}|p{5cm}|}
\hline
术语 & 英文 & 定义 \\
\hline
文档计划 & Documentation Plan & 规定项目文档编制范围、责任、进度和标准的计划文档 \\
文档编制 & Documentation Development & 按照标准要求编写、评审和发布文档的过程 \\
... & ... & ... \\
\hline
\end{longtable}

\hrule
\subsection{文档审批}
\begin{longtable}{|p{3cm}|p{4cm}|p{4cm}|p{3cm}|}
\hline
角色 & 姓名 & 签名 & 日期 \\
\hline
编制人 & [姓名] & [签名] & YYYY-MM-DD \\
审核人 & [姓名] & [签名] & YYYY-MM-DD \\
批准人 & [姓名] & [签名] & YYYY-MM-DD \\
\hline
\end{longtable}

\hrule
\textbf{文档状态}: 正式发布  
\textbf{下次审查日期}: YYYY-MM-DD  
\textbf{联系方式}: [邮箱/电话]


\end{document}
